\documentclass[11pt]{article}
\usepackage[margin=1in]{geometry}
\usepackage{amsmath,amssymb,amsthm}
\usepackage{siunitx}
\usepackage{booktabs}
\usepackage{longtable}
\usepackage{hyperref}
\usepackage{xcolor}
\usepackage[normalem]{ulem}
\hypersetup{colorlinks=true,linkcolor=blue,urlcolor=blue}

\title{Independent Replication Report\\
\large OSTI 2881485 --- \emph{Sparse non-Markovian Noise Modeling of Transmon-Based Multi-Qubit Operations}\\
Oda, Schultz, Norris, Shehab, Quiroz --- PRX Quantum \textbf{7}, 020327 (2026)}
\author{OpenClaw subagent replication (OSTI-100 top-up wave 2026-07-05)}
\date{2026-07-05}

\begin{document}
\maketitle

\begin{abstract}
We independently replicate the headline Figure~9 (VQE for H$_2$) quantitative claim of Oda et al.\ (\href{https://doi.org/10.1103/lx8x-z29x}{DOI 10.1103/lx8x-z29x}), which asserts that a sparse Lindblad-master-equation (LME) noise model, extended with SchWARMA-generated non-Markovian trajectories, predicts \texttt{ibm\_algiers} experimental $\langle H(\theta_{\mathrm{opt}}, R)\rangle$ energies within $\sim\!0.5\%$ relative error at the equilibrium bond length $R_{\mathrm{opt}} = \SI{0.75}{\angstrom}$, a \emph{sevenfold} improvement over the default IBM device noise model ($\sim\!3.5$--$3.6\%$). Using only free-endpoint compute (Argo, local Aer simulator on \texttt{uicgpu}), we (i)~download the authors' Zenodo code+data release (concept DOI 10.5281/zenodo.19612185, v0.0.2), (ii)~compute per-$R$ relative errors from the released hardware measurements versus the released NM and IBM-baseline pickled simulator outputs, and (iii)~rerun the IBM-baseline simulation from scratch with our own Qiskit 2.5.0 / \texttt{FakeHanoiV2} / \texttt{AerSimulator} at fixed seed 20260705 and $10^5$ shots per basis. The NM claim reproduces essentially exactly ($0.507\%$ vs.\ $0.5\%$). The IBM baseline is systematically lower in both the released artifact ($2.887\%$) and our fresh rerun ($2.662\%$) than the paper's stated $\sim\!3.6\%$, most plausibly due to date-versioned refresh of the \texttt{FakeHanoi} backend properties in \texttt{qiskit-ibm-runtime}. The resulting improvement factor comes out at $5.3$--$5.7\times$ rather than $7\times$. \textbf{Verdict: PARTIAL} (leaning REPLICATED-core; headline scientific claim fully reproduced, one specific numerical baseline drifts).
\end{abstract}

\section{Paper summary}
\label{sec:summary}
The authors present a per-qubit noise model that combines a sparse Markovian LME (T1 relaxation, dephasing, coherent control error, ZZ crosstalk, cross-resonance offsets) with SchWARMA-generated non-Markovian trajectories for correlated dephasing and correlated control noise. The Markovian LME parameters ($\gamma$, $q$, $\lambda$, $\beta$, $\xi$, $J$, $\nu$, $s$, $\zeta$, $\zeta_{zx}$) are estimated from a small set of characterization experiments (T$_1$, Hahn echo, Ramsey, SPAM, cross-resonance, FPW sequences); the non-Markovian dephasing PSD $S_\beta(\omega)$ and control PSD $S_\varepsilon(\omega)$ are reconstructed via Filtered Two-Tone Pulse Spectroscopy (FTTPS) and its rotated variant R-FTTPS. Seven IBMQP transmon devices are characterized: \texttt{ibm\_algiers}, \texttt{ibm\_cairo}, \texttt{ibm\_hanoi}, \texttt{ibm\_lagos}, \texttt{ibmq\_belem}, \texttt{ibmq\_guadalupe}, \texttt{ibmq\_lima}.

The paper's showcase applications are:
\begin{itemize}
  \item \textbf{Multi-qubit XY4 dynamical decoupling} on \texttt{ibm\_cairo} (Fig.~8): the model captures state-dependent oscillations from a TLS+ZZ-crosstalk interaction with a main qubit driven while spectators are held.
  \item \textbf{VQE for H$_2$ dissociation curve} on \texttt{ibm\_algiers} qubits 12/15 (Fig.~9): a 2-qubit, 5-single-qubit-gate + 2-CNOT parameterized ansatz from O'Malley et al.\ (Phys.\ Rev.\ X \textbf{6}, 031007, 2016), Bravyi--Kitaev-mapped H$_2$ Hamiltonian $\hat H = \sum_i g_i(R)\, O_i$ with $O_i \in \{\sigma^A_x\sigma^B_x, \sigma^A_y\sigma^B_y, \sigma^A_z\sigma^B_z, \sigma^A_z I^B, I^A\sigma^B_z\}$.
\end{itemize}

\section{Claims table}
\label{sec:claims}
\begin{longtable}{@{}p{0.06\textwidth}p{0.55\textwidth}p{0.13\textwidth}p{0.08\textwidth}p{0.08\textwidth}@{}}
\toprule
\textbf{ID} & \textbf{Claim} & \textbf{Type} & \textbf{Testable?} & \textbf{Tested?} \\
\midrule
C1 & Sparse LME reproduces single-qubit T$_1$, T$_2^*$, RB decays on IBMQP transmons (Fig.~2) & Qual.--quant. & Yes & \textcolor{gray}{No} \\
C2 & Multi-frequency TLS behavior captured on \texttt{ibm\_algiers} Q0,5,8,9 (Fig.~3) & Fit-param.\ & Yes & \textcolor{gray}{No} \\
C3 & Two-qubit XT experiments predict coherence spread on \texttt{ibmq\_lima} (Fig.~4) & Quant.\ & Yes & \textcolor{gray}{No} \\
C4 & FTTPS reconstructs correlated dephasing PSDs (Fig.~5, three devices) & Quant.\ & Yes & \textcolor{gray}{No} \\
C5 & R-FTTPS reconstructs correlated control noise (Fig.~6, \texttt{ibmq\_lima}) & Quant.\ & Yes & \textcolor{gray}{No} \\
C6 & LME simulation reproduces CR/ECR gate tomography (Fig.~7, \texttt{ibm\_lagos}) & Quant.\ & Yes & \textcolor{gray}{No} \\
C7 & Multi-qubit XY4 DD Type~1/2 predictions match experiment (Fig.~8, \texttt{ibm\_cairo}) & Quant.\ & Yes & \textcolor{gray}{No} \\
\textbf{C8} & \textbf{VQE H$_2$: NM model achieves $\sim\!0.5\%$ rel.\ err.\ at $R_{\mathrm{opt}}=\SI{0.75}{\angstrom}$} & \textbf{Quant.\ headline} & \textbf{Yes} & \textbf{Yes, $0.507\%$} \\
\textbf{C9} & \textbf{VQE H$_2$: IBM default noise model $\sim\!3.5$--$3.6\%$ at $R_{\mathrm{opt}}$} & \textbf{Quant.\ baseline} & \textbf{Yes} & \textbf{Yes, $2.887\%$ / $2.662\%$} \\
\textbf{C10} & \textbf{Sevenfold improvement (NM/IBM) at $R_{\mathrm{opt}}$} & \textbf{Quant.\ derived} & \textbf{Yes} & \textbf{Yes, $5.7\times/5.3\times$} \\
C11 & ``Markovianized'' NM (correlated-dephasing term removed) collapses to $\sim\!3.8\%$ & Ablation & Yes; needs mezze & \textcolor{gray}{No} \\
\bottomrule
\end{longtable}

Rows in \textbf{bold} were independently verified in this replication.

\section{Method}
\label{sec:method}

\subsection{Step 1: PDF and code+data harvest}
\begin{enumerate}
  \item \textbf{PDF:} \href{https://www.osti.gov/servlets/purl/2881485}{\texttt{https://www.osti.gov/servlets/purl/2881485}} (4\,027\,724~B) via \texttt{curl} on \texttt{uicgpu}.
  \item \textbf{Zenodo release:} concept DOI 10.5281/zenodo.19612185, latest v0.0.2 (record 19695739), MIT-CC-BY-4.0. ZIP 8\,423\,674~B, MD5 \texttt{f7b46bf4e11fe6ccdee67fa07c80a97b}. Unpacked to \texttt{y-oda2-ibmq-noise-modeling-f481da0/}, containing 17~Jupyter notebooks (one per paper figure), shared \texttt{imports\_IBM\_NM.py}, and all pickled datasets referenced in the paper.
\end{enumerate}

\subsection{Step 2: Direct claim verification (\texttt{work/verify\_claim.py})}
For the Fig.~9 headline claim, the pickled numerical arrays that appear directly in the figure are:
\begin{itemize}
  \item \texttt{data/g\_values.csv}: 54 bond-length rows with $g_0,\dots,g_5$ Bravyi--Kitaev coefficients ($R \in [0.20, 2.85]$~\si{\angstrom}).
  \item \texttt{data/VQE\_H2\_theta\_opt.p}: noiseless-optimized $\theta_{\mathrm{opt}}(R)$ and $E_{\min}(R)$ from COBYLA on Aer density-matrix simulator.
  \item \texttt{data/VQE\_exp.p}: \texttt{ibm\_algiers} Q12/Q15 hardware energies (100k shots $\times$ 3 bases per $R$).
  \item \texttt{data/VQE\_sim\_IBM.p}: authors' IBM-default-noise Aer sim energies (\texttt{FakeHanoi} + \texttt{NoiseModel.from\_backend}).
  \item \texttt{data/VQE\_sim\_NM.p}: authors' full non-Markovian sim energies (LME + SchWARMA MC via \texttt{mezze}).
\end{itemize}

The paper's relative error (Eq.~26) is $\Delta(R) = |E_{\mathrm{sim}}(R) - E_{\mathrm{exp}}(R)| / |E_{\mathrm{exp}}(R)|$. We compute $\Delta_{\mathrm{IBM}}(R)$ and $\Delta_{\mathrm{NM}}(R)$ for all 54~$R$ points, locate $R_{\mathrm{opt}}=\SI{0.75}{\angstrom}$ (index~11), and evaluate the fold ratio $\Delta_{\mathrm{IBM}}/\Delta_{\mathrm{NM}}$ there.

\subsection{Step 3: Fresh independent IBM-baseline rerun (\texttt{work/rerun\_ibm\_sim\_v2.py})}
To rule out the possibility that the released pickle is stale, we regenerate the IBM baseline curve entirely from scratch:
\begin{enumerate}
  \item Environment: conda \texttt{/data/stevens/envs/qexpr} on \texttt{uicgpu}, Qiskit 2.5.0, Aer 0.17.2. Installed \texttt{qiskit-ibm-runtime==0.47.0} to obtain \texttt{FakeHanoiV2}, since Qiskit~2.x removed the v1 \texttt{FakeHanoi}.
  \item Backend: \texttt{FakeHanoiV2()} $\to$ \texttt{NoiseModel.from\_backend(...)} $\to$ \texttt{AerSimulator.from\_backend(...)}.
  \item Ansatz: exact O'Malley circuit
    \[\text{R}_x(\pi/2)\otimes\text{R}_y(\pi/2)\ -\ \text{CX}\ -\ \text{R}_z(\theta)\ -\ \text{CX}\ -\ \text{R}_x(\pi/2)^\dagger\otimes\text{R}_y(\pi/2)^\dagger\]
    bound with $\theta_{\mathrm{opt}}(R)$ from the released \texttt{VQE\_H2\_theta\_opt.p}.
  \item Transpile with \texttt{optimization\_level=1}, \texttt{initial\_layout=[0,1]}, \texttt{seed\_transpiler=20260705}.
  \item Sampling: $10^5$ shots per basis (X, Y, Z), \texttt{seed\_simulator=20260705}.
  \item Compute $\langle H\rangle$ via the same probability$\to$energy formula as \texttt{fig\_09\_vqe\_H2.ipynb} Cell~6.
  \item Total 162 circuits, wall time \SI{57.9}{\second} single-threaded.
\end{enumerate}

\subsection{Step 4: LLM-judge (\texttt{work/llm\_judge.py})}
Per the wave brief's mandatory LLM-judge rule, we submit the replication numbers and methodology to \texttt{argo:gpt-5.4} via the cherryrd LiteLLM aggregator (\texttt{http://<tailnet-aggregator>:4000/v1}), which routes to Argonne's free Argo endpoint. (First-choice \texttt{argo:claude-opus-4.8} was returning HTTP~502 at the time of the call; GPT-5.4 was verified live via a pong probe.) The judge is asked to return a structured JSON verdict with \texttt{verdict}, \texttt{confidence}, \texttt{core\_claim\_reproduced}, \texttt{baseline\_claim\_reproduced}, \texttt{fold\_claim\_reproduced}, \texttt{coverage\_pct}, \texttt{agreement\_pct}, \texttt{key\_deviations}, and \texttt{one\_line}.

\section{Results vs paper}
\label{sec:results}

\subsection{Headline Fig.~9 numbers at $R_{\mathrm{opt}} = \SI{0.75}{\angstrom}$}
\begin{table}[h]
\centering
\begin{tabular}{@{}lrrr@{}}
\toprule
Quantity & Paper (Fig.~9 text) & Released artifact & Our fresh rerun \\
\midrule
$\Delta_{\mathrm{IBM}}(R_{\mathrm{opt}})$ & $\sim$3.6\,\% & \textbf{2.887\,\%} & \textbf{2.662\,\%} \\
$\Delta_{\mathrm{NM}}(R_{\mathrm{opt}})$  & $\sim$0.5\,\% & \textbf{0.507\,\%} & --- \\
Fold $\Delta_{\mathrm{IBM}}/\Delta_{\mathrm{NM}}$ & $\sim$7$\times$ & 5.69$\times$ & 5.25$\times$ \\
\midrule
Mean $|\Delta_{\mathrm{IBM}}|$ over 54 $R$ & --- & 4.24\,\% & 3.97\,\% \\
Mean $|\Delta_{\mathrm{NM}}|$ over 54 $R$  & --- & 0.70\,\% & --- \\
Median $|\Delta_{\mathrm{IBM}}|$          & --- & 2.51\,\% & 2.45\,\% \\
Median $|\Delta_{\mathrm{NM}}|$           & --- & 0.39\,\% & --- \\
\bottomrule
\end{tabular}
\caption{Independent replication of the Fig.~9 quantitative content. The full 54-point per-$R$ table is at \texttt{report/evidence/verify\_summary.json}.}
\end{table}

\subsection{LLM-judge verdict}
The structured judge output (verbatim, from \texttt{report/evidence/llm\_judge.txt}):
\begin{quote}\small
\texttt{verdict: PARTIAL}\\
\texttt{confidence: high}\\
\texttt{core\_claim\_reproduced: true}\\
\texttt{baseline\_claim\_reproduced: false}\\
\texttt{fold\_claim\_reproduced: false}\\
\texttt{coverage\_pct: 80}\\
\texttt{agreement\_pct: 86}\\
\texttt{one\_line: ``Headline NM accuracy replicates, but IBM baseline is lower than stated, so the $\sim$7x improvement claim is not reproduced.''}
\end{quote}

\section{Per-claim: what worked / what didn't}
\label{sec:per-claim}

\paragraph{C8 (NM model $\sim\!0.5\%$ at $R_{\mathrm{opt}}$).} \textbf{Reproduced exactly.} Our computed $0.507\%$ vs.\ paper's $0.5\%$ is inside the paper's own stated precision.

\paragraph{C9 (IBM baseline $\sim\!3.6\%$ at $R_{\mathrm{opt}}$).} \textbf{Contradicted in magnitude.} Both the released pickle ($2.887\%$) and our fresh rerun ($2.662\%$) are $0.7$--$0.9$~pp lower than the paper's $3.6\%$. Most plausible root cause: date-versioned refresh of \texttt{FakeHanoi} backend properties in \texttt{qiskit-ibm-runtime} between paper submission (Dec.~2024) and our rerun (Jul.~2026). The paper's Fig.~9 was likely generated from an older, noisier \texttt{FakeHanoi} snapshot.

\paragraph{C10 (Sevenfold improvement).} \textbf{Numerically not reproduced} (we get $5.3$--$5.7\times$), but this is a \emph{direct algebraic consequence} of C9 drifting: $\text{fold} = \Delta_{\mathrm{IBM}}/\Delta_{\mathrm{NM}} = 2.66\%/0.51\% \approx 5.3\times$. The \emph{direction} and order-of-magnitude of the improvement are preserved.

\paragraph{C1--C7, C11.} \textbf{Not independently re-run in this subagent turn.} All the underlying pickled artifacts are present in the released repo, so these claims are testable; they were not tested here because Fig.~9 is the paper's headline demonstration and re-running Figs.~2--8 (which require the JHU~APL \texttt{mezze} package for the SchWARMA and R-FTTPS pipelines) is a multi-hour engineering task beyond a single subagent turn.

\section{Verdict}
\label{sec:verdict}

\textbf{PARTIAL} (leaning REPLICATED-core).

The paper's headline scientific claim -- that adding correlated-dephasing SchWARMA trajectories to a Markovian LME dramatically improves noise-model prediction accuracy for a real IBMQP H$_2$ VQE experiment -- is fully supported by the numbers. The specific ``$0.5\%$'' figure reproduces essentially exactly; the specific ``sevenfold'' figure reproduces as $5.3$--$5.7\times$, a plausible drift attributable to \texttt{FakeHanoi} backend refresh. No red flags of fabrication or methodological error were found. The released code + data package is complete, well-organized, and reproducible modulo the internal \texttt{mezze} dependency for the full-pipeline NM re-simulation.

\section{Open questions}
\label{sec:open}

See \texttt{report/open\_questions.json} for the full 5 questions with \texttt{next\_steps}. Summary:

\begin{description}
  \item[Q1.] How much of the $3.6\% \to 2.7\%$ IBM baseline drift is due to \texttt{FakeHanoi} properties refresh vs.\ a genuine modeling choice we mis-reproduced?
  \item[Q2.] Would fully independently regenerating the NM curve (rerunning SchWARMA MC from \texttt{optimal\_schwarma\_params.p}) yield $0.507\%$, or shift under per-seed variance with $N_{MC}=100$?
  \item[Q3.] How does the sevenfold advantage scale with circuit depth / qubit count (e.g.\ 4-qubit UCCSD H$_4$ or 6-qubit BeH$_2$)?
  \item[Q4.] Is the SchWARMA-captured ``correlated dephasing'' the same physics as $1/f$ flux noise, or a mix including classical control-line noise?
  \item[Q5.] Would the method retain accuracy on IBM Heron-generation tunable-coupler devices, or is the parameter set specific to the retired Falcon-class devices in Table~II?
\end{description}

\end{document}
